% Section 2: Related Work
% Source: context/Related Work Literature Review.md (all 5 areas)
% Source: context/ICSA_RA_Literature_Review.md §11 (reference quality, paper-section mapping)
% Source: context/ICSA_RA_Literature_Review.md §7.3 (Swiss Cheese differentiation)
% Source: context/ICSA_RA_Literature_Review.md §7.4 (ACE citation strategy)
% Source: context/ICSA_RA_Literature_Review.md §11.4 (Chen et al., Abadeh et al. differentiation)
%
% EM-DASH PROHIBITION: Do not use --- or \textemdash. Use colons or semicolons.

\section{Related Work}\label{sec:related_work}

\subsubsection*{LLM Output Fidelity}
LLM outputs are unreliable in ways that extend beyond factual
hallucination~\cite{huang2025hallucination, ji2023hallucination}. Nie et
al.~\cite{nie2025decoding} study how code LLMs handle secrets at the token
level, focusing on memorisation risk: the unintended reproduction of
training-data credentials. No prior work studies the complementary failure: the
\emph{corruption} of novel, unseen credentials when an LLM relays them at
inference time through a tool chain. This gap defines the empirical contribution
of the present study.

\subsubsection*{Tokenisation and String Integrity}
Byte Pair Encoding (BPE)~\cite{sennrich2016bpe} segments input into subword
tokens using merge rules from a training corpus. Applied to high-entropy strings
such as JWTs, BPE produces a non-uniform distribution that Chen et
al.~\cite{chen2026gibberish} term \emph{gibberish bias}: accidental
natural-language fragments (e.g., ``cut'' inside Base64url) merge into single
tokens while adjacent characters receive individual slots. Chen et al.\ measure
KL divergence of 2.668 between credential-string and normal-code token
distributions, and show the effect intensifies with vocabulary size (the current
scaling trajectory). We hypothesise this same distributional shift underlies
the relay fidelity failures we observe: rather than enabling memorised secret
reproduction, it impedes accurate relay of novel credential strings.
Chen et al.'s work concerns training-data secret \emph{leakage}, not
inference-time relay of novel credentials.

\subsubsection*{MCP and Agentic AI Security}
The Model Context Protocol~\cite{anthropic2024mcp} enables LLMs to invoke
external tools through structured function calls. Schick et
al.~\cite{schick2023toolformer} established that LLMs can learn tool use; He et
al.~\cite{he2025security} survey the resulting security and privacy risks
across the broader LLM agent landscape. Abadeh et al.~\cite{abadeh2026mcp}
provide the first peer-reviewed study of MCP-specific vulnerabilities, scanning
server configurations for privilege escalation and injection risks. Their
approach is complementary: static scanning examines infrastructure-level
configuration; our study addresses a runtime failure mode that arises during
live LLM inference and falls outside the scope of static analysis. Greshake et
al.~\cite{greshake2023prompt} and Liu et al.~\cite{liu2024prompt} study prompt
injection attacks on LLM-integrated applications; the Vault architecture we
propose incidentally mitigates indirect prompt injection targeting credentials,
as the credential never enters the LLM context.

\subsubsection*{Function-Call and Tool-Use Reliability}
Song et al.~\cite{song2025callnavi} find systematic argument misgeneration
in LLM function-calling across diverse API schemas; Yu
et al.~\cite{yu2026wildtoolbench} report compounding errors in multi-step
tool-use, with both studies tolerating near-miss outputs that constitute a hard
cryptographic failure in our setting. Haque et al.~\cite{haque2026verbatim}
study high-entropy string transcription failures in LLM code generation; that
work does not address credential relay.

\subsubsection*{Credential Relay in Distributed Systems}
The challenge of securely propagating bearer tokens across distributed
components is well-documented in microservice architectures. Yarygina and
Bagge~\cite{yarygina2018microservice} identify token propagation as a primary
security challenge in microservices. These patterns exist for deterministic
microservices, but they fail when a probabilistic LLM is introduced into the
relay chain: the LLM is not a transparent relay component but an active
transformer of its input. To our knowledge, no prior work addresses bearer
token integrity when the relay component is a stochastic language model.

\subsubsection*{Reference Architectures for LLM Agent Systems}
The Swiss Cheese Model for AI Safety~\cite{shamsujjoha2025swisscheese},
published at ICSA~2025, proposes a multi-layer guardrail reference architecture
for foundation model agents, inspired by Reason's accident causation model~\cite{reason1990human}. While
conceptually adjacent, it addresses general agent failure modes (hallucination,
misalignment, prompt injection) through stacked defensive layers. The failure
mode we study is distinct: verbatim credential relay corruption produces
well-formed but cryptographically invalid output that passes format validation at
every Swiss Cheese layer. Our architectural response is also distinct: rather
than adding layers around LLM inference, the Credential Interception Vault
removes the credential from the inference path entirely. The two architectures
are complementary. Li et al.'s ACE architecture~\cite{li2026ace}, published at
NDSS~2026, isolates LLM plan generation from deterministic execution to prevent
prompt injection. We apply the same separation principle to a different failure
mode: verbatim string relay fidelity. In both cases, deterministic correctness
guarantees cannot be provided by a probabilistic component, regardless of the
operation's apparent simplicity.

Grammar-constrained decoding~\cite{geng2023grammar} cannot guarantee byte-level
fidelity of opaque strings: BPE tokenisation fragments the credential before
generation begins, and no grammar can reconstruct a specific byte sequence
already split across merge boundaries. OAuth~2.0 Token Exchange
(RFC~8693~\cite{rfc8693}) assumes a deterministic relay agent; the stochastic
byte-level transformation we measure is outside its threat model.

The RA canon provides the classification framework for our contribution.
Following Angelov et al.~\cite{angelov2012framework}, we classify the Vault as
a Type~2 facilitation-type RA: designed by researchers to guide practitioners,
with voluntary adoption. Galster and Avgeriou~\cite{galster2011empirically}
propose a six-step methodology for empirically-grounded RAs; our study
satisfies all six steps (Section~\ref{sec:ra_type}).


